\documentclass[11pt]{amsart}
\usepackage[margin=1.1in]{geometry}
\usepackage{amsmath,amssymb,amsthm,mathtools}
\usepackage[hidelinks]{hyperref}
\usepackage{orcidlink}

\newtheorem{theorem}{Theorem}[section]
\newtheorem{lemma}[theorem]{Lemma}
\theoremstyle{definition}
\newtheorem{definition}[theorem]{Definition}
\theoremstyle{remark}
\newtheorem{remark}[theorem]{Remark}

\title[Essential vectors and rank]{Essential vectors and rank:\\ a machine-checked note}
\author{Ella Agnes Puerto\,\orcidlink{0009-0005-1624-2025}}
\address{Independent researcher}
\email{eagnespuerto@gmail.com}
\date{November 2026. Source, data and this note: \href{https://doi.org/10.5281/zenodo.24475596}{doi:10.5281/zenodo.24475596}, released under the MIT license}

\begin{document}

\begin{abstract}
Take a finite family of nonzero vectors over a field, and suppose at least one is redundant,
that is, it lies in the span of the others. Call a vector \emph{essential} when the others
cannot reproduce it. I show that the number of essential vectors, plus one, is at most the
rank of the family. In matroid language the essential vectors are the coloops, and the
statement is that a representable matroid with a circuit has strictly fewer coloops than its
rank. The proof is formalized in Lean~4 against Mathlib and uses only the axioms
\texttt{propext}, \texttt{Classical.choice} and \texttt{Quot.sound}.
\end{abstract}

\maketitle

\section{Introduction}\label{sec:intro}
The fact is elementary, so I lead with the picture.

Think of the vectors as people on a small crew, each able to pull in some direction of work. A
person is \emph{redundant} when the rest of the crew together already covers everything they
do. A person is \emph{essential} when no mix of the others can stand in for them. The essential
people are the ones I cannot drop without losing reach. Two things follow. No essential person
wastes a slot, so each one adds a direction the rest did not have; that is why their count is
at most the rank. And once some person is redundant, the crew has slack: the essential people
alone do not use up all the room, so the reach beats their count by at least one.

The same idea works with roads. Each vector is a road out of town, and the rank is how many
different places I can reach by combining roads. An essential road goes somewhere the others
cannot. If one road is redundant, meaning I could already reach its end another way, then the
roads together reach strictly more places than there are essential roads.

\smallskip\noindent\textbf{Relation to matroid theory.}
An essential vector is a \emph{coloop} of the matroid that the family represents: an element
lying in every basis, equivalently a vector outside the closure (here, the span) of the
remaining elements~\cite{Oxley}. The coloops of a matroid form an independent set, and a
matroid is free exactly when every element is a coloop. Theorem~\ref{thm:main} is the
quantitative form: once the family carries a circuit (a redundant vector), its coloops fall
short of a basis, so their number drops below the rank. The argument is standard; the
contribution is a short, fully checked proof in a shape that drops into a larger formal
development.

\section{Statement}\label{sec:statement}
Let $K$ be a field and $V$ a vector space over $K$. Let $r\colon\iota\to V$ be a family of
vectors indexed by a finite type $\iota$. Assume $r_i\neq 0$ for every $i$, and assume one
vector is redundant: some $r_j$ lies in $\operatorname{span}_K\{\,r_\ell:\ell\neq j\,\}$.

\begin{definition}\label{def:essential}
An index $i$ is \emph{essential} when $r_i\notin\operatorname{span}_K\{\,r_\ell:\ell\neq i\,\}$.
\end{definition}

Write $e$ for the number of essential indices and $V_0=\operatorname{span}_K\{\,r_i:i\in\iota\,\}$
for the span of the whole family. A finite set spans $V_0$, so it is finite-dimensional; no
assumption on the ambient space $V$ is needed.

\begin{theorem}\label{thm:main}
$e+1\le\dim_K V_0$.
\end{theorem}

\section{Proof}\label{sec:proof}
\begin{proof}
\textbf{The essential vectors are independent.} Suppose a combination of the essential vectors
vanishes,
\[ \sum_t c_t\, r_t = 0, \qquad t \text{ over the essential indices.} \]
If some coefficient $c_{t_0}$ is nonzero, solve for $r_{t_0}$:
\[ r_{t_0} = -\,c_{t_0}^{-1} \sum_{t\neq t_0} c_t\, r_t. \]
Every term on the right is a scalar times $r_t$ with $t\neq t_0$, so
$r_{t_0}\in\operatorname{span}_K\{\,r_\ell:\ell\neq t_0\,\}$, contradicting that $t_0$ is
essential. So all coefficients vanish, the essential vectors are linearly independent,
$\dim_K(\text{span of the essentials})=e$, and that span lies in $V_0$, giving $e\le\dim_K V_0$.

\smallskip
\textbf{The essentials do not span.} Let $r_j$ be the redundant vector from the hypothesis; it
is nonzero. Suppose $r_j$ were in the span of the essentials,
\[ r_j = \sum_t c_t\, r_t, \qquad t \text{ over the essential indices.} \]
Not every $c_t$ is zero, since $r_j\neq 0$. Pick $t_0$ with $c_{t_0}\neq 0$ and solve:
\[ r_{t_0} = c_{t_0}^{-1}\Bigl( r_j - \sum_{t\neq t_0} c_t\, r_t \Bigr). \]
The index $j$ is not essential, so $j\neq t_0$, and every index on the right differs from
$t_0$; hence $r_{t_0}\in\operatorname{span}_K\{\,r_\ell:\ell\neq t_0\,\}$, again contradicting
that $t_0$ is essential. So $r_j$ lies in $V_0$ but outside the span of the essentials, making
that span a proper subspace of $V_0$. A strict inclusion of finite-dimensional subspaces loses
at least one dimension, so $e=\dim_K(\text{span of the essentials})<\dim_K V_0$, that is,
$e+1\le\dim_K V_0$.
\end{proof}

\section{Formalization}\label{sec:formal}
I formalized Theorem~\ref{thm:main} in Lean~4~\cite{Lean4} with Mathlib~\cite{Mathlib} as the
theorem \texttt{card\_essential\_add\_one\_le\_finrank\_span}. The two steps reuse one move:
solve a single vector out of a dependence and show it lands in the span of the family with that
index removed. In Lean that is \texttt{Fintype.linearIndependent\_iff} for the first step and
\texttt{mem\_span\_range\_iff\_exists\_fun} for the second, each followed by
\texttt{Finset.sum\_erase\_add} to split off the chosen term. The conclusion uses
\texttt{finrank\_span\_eq\_card}, \texttt{Submodule.finrank\_lt\_finrank\_of\_lt} and
\texttt{FiniteDimensional.span\_of\_finite}.

Mathlib also carries a matroid library, developed by Nelson and contributors~\cite{Nelson}, in
which the statement could be phrased with coloops; I kept it in direct linear-algebra terms so
that it imports into any development that already speaks of \texttt{finrank}.

\smallskip\noindent\textbf{Provenance.}
This lemma began as an open gap in the public pequalnp.com collaboration~\cite{pequalnp}, an
attempt at P versus NP, where it was stated over the rationals. The thread for that gap was
closed before the proof could be merged, so I generalized it to an arbitrary field and record
it here on its own.

\section*{Acknowledgements and use of AI}
This note and its Lean proof were prepared by the author with extensive use of the AI assistant
Claude (Anthropic; model Claude Opus~4.8), through Claude Code: the formalization, the
write-up, and the checking of references against primary sources.

\appendix
\section{Reproducibility}\label{app:repro}
The Lean source, this note and a build configuration are in the repository~\cite{Repo} and
archived on Zenodo~\cite{Zenodo}. The proof was checked with Lean~\texttt{v4.36.0} and Mathlib
pinned to commit \texttt{fc4f59e} (the v4.36.0 release). It has no \texttt{sorry} and no
\texttt{native\_decide}; \texttt{\#print axioms} reports only \texttt{propext},
\texttt{Classical.choice} and \texttt{Quot.sound}. To rebuild, run \texttt{lake exe cache get}
and then \texttt{lake build}.

\begin{thebibliography}{9}
\bibitem{Lean4} L.~de~Moura and S.~Ullrich, \emph{The Lean~4 theorem prover and programming
  language}, in Automated Deduction---CADE~28 (A.~Platzer and G.~Sutcliffe, eds.), Lecture
  Notes in Comput.\ Sci., Springer, 2021, pp.~625--635;
  \href{https://doi.org/10.1007/978-3-030-79876-5_37}{doi:10.1007/978-3-030-79876-5\_37}.
\bibitem{Mathlib} The mathlib Community, \emph{The Lean mathematical library}, in Proc.\ 9th
  ACM SIGPLAN Int.\ Conf.\ on Certified Programs and Proofs (CPP~2020), 2020, pp.~367--381;
  \href{https://doi.org/10.1145/3372885.3373824}{doi:10.1145/3372885.3373824}; arXiv:1910.09336.
\bibitem{Oxley} J.~G.~Oxley, \emph{Matroid Theory}, 2nd ed., Oxford Graduate Texts in Math.\
  \textbf{21}, Oxford Univ.\ Press, 2011.
\bibitem{Nelson} P.~Nelson and contributors, \emph{A formalization of matroids in Lean},
  development repository, \url{https://github.com/apnelson1/Matroid}; much of it integrated
  into Mathlib~\cite{Mathlib}.
\bibitem{pequalnp} The pequalnp collaboration, \emph{A public attempt at P versus NP},
  \url{https://pequalnp.com}; proofs repository \url{https://github.com/brilliant-hq/pequalnp-proofs}.
\bibitem{Repo} E.~A.~Puerto, \emph{Essential vectors and rank} (source repository),
  \url{https://github.com/eagnespuerto/essential-vectors-rank}.
\bibitem{Zenodo} E.~A.~Puerto, \emph{Essential vectors and rank: a machine-checked Lean~4
  proof}, Zenodo, 2026; \href{https://doi.org/10.5281/zenodo.24475596}{doi:10.5281/zenodo.24475596}.
\end{thebibliography}

\end{document}
